% Based on the typography, KOMA-Script class, colors, headers and 6x9 dimensions
% of Non-Deterministic Loop Engineering. Rebuilt to be self-contained.
\usepackage[english]{babel}
% Latin Modern Sans is the book's original face. Use its Unicode OpenType version
% with XeLaTeX so euros, arrows and typographic punctuation are not silently lost.
\usepackage{fontspec}
\setmainfont{lmsans10-regular.otf}[BoldFont=lmsans10-bold.otf, ItalicFont=lmsans10-oblique.otf, BoldItalicFont=lmsans10-boldoblique.otf]
\setsansfont{lmsans10-regular.otf}[BoldFont=lmsans10-bold.otf, ItalicFont=lmsans10-oblique.otf, BoldItalicFont=lmsans10-boldoblique.otf]
\setmonofont{lmmono10-regular.otf}[BoldFont=lmmonolt10-bold.otf, ItalicFont=lmmono10-italic.otf]
\renewcommand{\familydefault}{\sfdefault}
\usepackage{newunicodechar}
\newunicodechar{χ}{\ensuremath{\chi}}
\newunicodechar{≈}{\ensuremath{\approx}}
\newunicodechar{⁻}{\textsuperscript{-}}
\newunicodechar{⁶}{\textsuperscript{6}}
\usepackage{textcomp}
\usepackage[final]{microtype}
\usepackage[table]{xcolor}
\definecolor{primaryblue}{RGB}{41,128,185}
\definecolor{secondaryblue}{RGB}{52,152,219}
\definecolor{darkgray}{RGB}{52,73,94}
\definecolor{lightgray}{RGB}{236,240,241}
\definecolor{successgreen}{RGB}{39,174,96}
\definecolor{warningyellow}{RGB}{241,196,15}
\usepackage[
 paperwidth=\bookwidth,paperheight=\bookheight,
 top=0.75in,bottom=0.75in,outer=0.625in,inner=0.875in,
 includehead,includefoot,heightrounded
]{geometry}
\usepackage{setspace}
\setstretch{1.055}
\usepackage{ragged2e}
\AtBeginDocument{\justifying}
\usepackage[automark,headsepline]{scrlayer-scrpage}
\clearpairofpagestyles
\automark[chapter]{chapter}
\ohead{\footnotesize\itshape\leftmark}
\cfoot*{\pagemark}
\renewcommand*{\chapterpagestyle}{plain}
\renewcommand*{\partpagestyle}{empty}
\renewcommand*{\chaptermarkformat}{}
\setkomafont{pageheadfoot}{\footnotesize\normalfont\sffamily}
\setkomafont{pagenumber}{\small\normalfont\sffamily}
\setkomafont{disposition}{\normalfont\sffamily\bfseries}
\setkomafont{chapter}{\normalfont\sffamily\huge\bfseries\color{primaryblue}}
\setkomafont{section}{\normalfont\sffamily\Large\bfseries\color{darkgray}}
\setkomafont{subsection}{\normalfont\sffamily\large\bfseries\color{darkgray}}
\setkomafont{part}{\normalfont\sffamily\Huge\bfseries\color{primaryblue}}
\RedeclareSectionCommand[beforeskip=0pt,afterskip=2.2\baselineskip]{chapter}
\RedeclareSectionCommand[beforeskip=-0.8\baselineskip,afterskip=.45\baselineskip]{section}
\RedeclareSectionCommand[beforeskip=-0.6\baselineskip,afterskip=.25\baselineskip]{subsection}
\usepackage{amsmath,amssymb,amsthm}
\usepackage{graphicx,booktabs,tabularx,longtable,array,multirow}
\usepackage{enumitem}
\setlist{topsep=0.25em,partopsep=0pt,itemsep=0.15em,parsep=0pt,leftmargin=*}
\usepackage{fancyvrb}
\usepackage{fvextra}
\DefineVerbatimEnvironment{verbatim}{Verbatim}{breaklines=true,breakanywhere=true,fontsize=\footnotesize}
\usepackage{csquotes}
\usepackage{float} % stable placement for original practitioner agency schematic
\usepackage{eso-pic}
\usepackage[most]{tcolorbox}

% Didactic edition: explanatory boxes can split naturally across pages.
\newtcolorbox{conceptbox}[1]{enhanced,breakable,colback=primaryblue!4,
 colframe=primaryblue!60,colbacktitle=primaryblue!13,coltitle=darkgray,
 boxrule=.5pt,arc=.8mm,title={Concept: #1},fonttitle=\bfseries\small,
 before skip=8pt,after skip=8pt,left=1.7mm,right=1.7mm,top=1mm,bottom=1mm}
\newtcolorbox{examplebox}[1]{enhanced,breakable,colback=successgreen!3,
 colframe=successgreen!52!black,colbacktitle=successgreen!11,coltitle=darkgray,
 boxrule=.5pt,arc=.8mm,title={Worked example: #1},fonttitle=\bfseries\small,
 before skip=8pt,after skip=8pt,left=1.7mm,right=1.7mm,top=1mm,bottom=1mm}
\newtcolorbox{practicebox}[1]{enhanced,breakable,colback=warningyellow!6,
 colframe=warningyellow!60!black,colbacktitle=warningyellow!16,coltitle=darkgray,
 boxrule=.5pt,arc=.8mm,title={Check your understanding: #1},fonttitle=\bfseries\small,
 before skip=8pt,after skip=8pt,left=1.7mm,right=1.7mm,top=1mm,bottom=1mm}

\newtcolorbox{bookconcept}[1][]{colback=primaryblue!4,colframe=primaryblue!60,boxrule=0.5pt,title={#1},fonttitle=\bfseries}
% Short learning and key-concept boxes stay intact across page boundaries.
\newtcolorbox{bookkeyconcept}[1][]{%
 colback=lightgray!32,colframe=primaryblue!54,colbacktitle=primaryblue!11,
 coltitle=darkgray,boxrule=.48pt,title={#1},fonttitle=\bfseries\footnotesize,
 fontupper=\small,left=1.4mm,right=1.4mm,top=.65mm,bottom=.65mm,
 before skip=.65em,after skip=.65em}

\usepackage[backend=biber,style=numeric-comp,sorting=nyt,doi=true,url=true,isbn=true,maxnames=2,giveninits=true]{biblatex}
\addbibresource{bib/references.bib} % All fourteen chapters; aliases preserve checkpoint keys.
\newcommand{\oroszcite}[1]{\workref{#1}\,\autocite{#1}}

% Selected stable Atlas identifiers, separate from BibLaTeX's generated numbers.
% The source guide explains the scope; this book does not reproduce the full Atlas.
\DeclareFieldFormat{bookworkid}{\textcolor{primaryblue}{\sffamily\bfseries #1}}
\renewbibmacro*{begentry}{%
  \iffieldundef{usera}{}{% 
    \printfield[bookworkid]{usera}\addspace
    \textcolor{darkgray}{\textbullet}\addspace
  }%
}
% Print the stable identifier carried by a cited entry's usera field.
\DeclareCiteCommand{\workref}
  {}
  {\printtext[bibhyperref]{\printfield[bookworkid]{usera}}}
  {\multicitedelim}
  {}
% Cite Osmani's durable Oxxx ID beside the generated bibliography number.
\newcommand{\osmanicite}[1]{\workref{#1}\,\autocite{#1}}
\defbibnote{sourceidlegend}{%
  \noindent\small\textbf{Reading the references.}
  Bracketed numbers are generated for this edition. Blue identifiers, where supplied, locate a selected record in the Knowledge Atlas. They are source identifiers, not evidence ratings. Appendix~\ref{sec:source-guide} explains how to use sources and practice guides to investigate an engineering question.\par\medskip
}

% Resolve original editorial work-ID genitives as authored attribution with a citation,
% avoiding typographically and grammatically incorrect expressions such as `[38]'s`.
\DeclareRobustCommand{\workposs}[1]{\citeauthor{#1}\textquotesingle s~\autocite{#1}}
\usepackage{hyperref}
\hypersetup{colorlinks=true,linkcolor=primaryblue,citecolor=primaryblue,urlcolor=secondaryblue,pdftitle={\booktitle},pdfsubject={\booksubtitle},pdfauthor={\authorname},pdfkeywords={software engineering, AI coding agents, independent verification, accepted change}}
\usepackage{bookmark}
\urlstyle{same}
\tolerance=1300
\emergencystretch=2em
\hbadness=1900

% Reader exercise metadata; worked examples are placed in their chapters.
\newcommand{\exercisemeta}[2]{\noindent\textbf{Time:} #1\hspace{1.4em}\textbf{Deliverable:} #2\par\smallskip}

% Three-digit page numbers need a wider TOC field; retain the existing type size.
\makeatletter
\renewcommand{\@pnumwidth}{2.5em}
\renewcommand{\@tocrmarg}{3.5em}
\makeatother

% Prevent 10.10/11.10 and B.2.12 labels from touching TOC titles.
\DeclareTOCStyleEntry[numwidth=3.0em]{tocline}{section}
\DeclareTOCStyleEntry[numwidth=3.8em]{tocline}{subsection}


% --- Engineering Software with Agentic AI book adaptations -----------------
% NDSE2 base uses 6x9 KOMA, Latin Modern Sans, blue headings, headers and callouts.
% All fourteen chapters use these listings, diagrams and named boxes.
\usepackage{needspace}
\usepackage{xurl}
\usepackage{tikz}
\usetikzlibrary{arrows.meta,positioning,calc,fit}
\usepackage{listings}
\colorlet{ink}{darkgray}
\colorlet{accent}{primaryblue}
\colorlet{pale}{lightgray}
\definecolor{codebg}{HTML}{F3F5F5}
\lstdefinestyle{lab}{basicstyle=\ttfamily\scriptsize,breaklines=true,backgroundcolor=\color{codebg},frame=single,framerule=0.25pt,rulecolor=\color{darkgray!40},xleftmargin=4pt,xrightmargin=4pt,aboveskip=6pt,belowskip=6pt,columns=fullflexible,keepspaces=true,showstringspaces=false}
\lstset{style=lab}
% These short callouts each fit on one page; keep their argument intact.
\newtcolorbox{readerbox}[1]{enhanced,colback=primaryblue!4,colframe=primaryblue!58,boxrule=.5pt,arc=.8mm,title={#1},fonttitle=\bfseries\footnotesize,before skip=6pt,after skip=6pt}
\newtcolorbox{warningbox}[1]{enhanced,colback=warningyellow!8,colframe=warningyellow!70!black,boxrule=.5pt,arc=.8mm,title={#1},fonttitle=\bfseries\footnotesize,before skip=6pt,after skip=6pt}
\setcounter{tocdepth}{1}
\emergencystretch=3em
